\section*{\texorpdfstring{\S\CurrentExerciseGroup}{第{\CurrentExerciseGroup}节}}
\phantomsection
\addcontentsline{toc}{section}{第{\CurrentExerciseGroup}节习题}

\begin{enumerate}
\def\labelenumi{\arabic{enumi})}
\item
  设 T 和 X 是两个紧空间，\(\pi\) 是 T 到 X 中的一个连续映射，g 是定义在 T 上的连续且有限的数值函数。对每个 \(x \in X\) ，用 \(f(x)\) 表示 \(g(t)\) 在集 \(\overline{\pi}^{-1}(x)\) 上的下确界。给出一个 f 不连续的例子。（取 T = R 中的 {[}0,1{]}，对 \(\pi\) 取把 T 中 0 与 1 等同起来所得商空间上 T 的典范映射。）
\item
  设 X 和 \(\nu\) 是 Ch. IV, §1 习题 5 中定义的局部紧空间和测度；设 f 是 X 中由点 \((0,y)\) 组成的集 D 的特征函数。证明：在下列两种情形下，\(\nu=\int g(t)\varepsilon_{\pi(t)}d\mu(t)\) ，且
\end{enumerate}

\[
\int^ {*} f (\pi (t)) g (t) d \mu (t) <   \int^ {*} f (x) d \nu (x) = + \infty .
\]

\begin{enumerate}
\def\labelenumi{\alph{enumi})}
\tightlist
\item
  取 T = X，对 \(\mu\) 取由质量 \(\log n/n^{3}\) （在每点 \((1/n, k/n^{2})\) 处）所定义的测度。对 \(\pi\) 取恒等映射，g 由下列条件定义：
\end{enumerate}

\[
g (0, y) = 0, \quad g (1 / n, k / n ^ {2}) = 1 / \log n;
\]

\(g\) 是连续的，但 \(g(t) > 0\) 并非对每个 \(t \in \mathbf{T}\) 成立。

\begin{enumerate}
\def\labelenumi{\alph{enumi})}
\setcounter{enumi}{1}
\tightlist
\item
  取 T 为赋以离散拓扑的集 X，对 \(\mu\) 取由与 \(\nu\) 相同的质量定义的测度。取 \(g(t) = 1\) （对所有 \(t\) ），对 \(\pi\) 取恒等映射：后者是连续的，但不是逆紧的。
\end{enumerate}
% semantic-fragments: legacy-input-seam
\relax{}
\setcounter{exercisegroup}{5}
\edef\CurrentExerciseGroup{\arabic{exercisegroup}}
